% Канонический TikZ-рисунок варианта 03. Править только здесь.
\long\def\figStereoV03{%
\begin{tikzpicture}[
    scale=0.9,
    line cap=round,
    line join=round,
    visible/.style={
        line width=0.9pt
    },
    hidden/.style={
        line width=0.7pt,
        dashed,
        dash pattern=on 3pt off 2pt
    }
]

% -------------------------------------------------
% Вершины пирамиды
% -------------------------------------------------

\coordinate (S) at (2.10,4.70);

% Прямоугольное основание
\coordinate (A) at (0.55,0.90);
\coordinate (B) at (2.10,0.25);
\coordinate (C) at (4.55,1.35);
\coordinate (D) at (2.85,2.05);

% Основание высоты лежит на заднем ребре AD
\coordinate (H) at ($(A)!0.58!(D)$);

% -------------------------------------------------
% Скрытые рёбра основания
% -------------------------------------------------

\draw[hidden] (A)--(D);
\draw[hidden] (D)--(C);

% -------------------------------------------------
% Скрытое боковое ребро
% -------------------------------------------------

\draw[hidden] (S)--(D);

% -------------------------------------------------
% Высота пирамиды
% Она лежит в боковой грани SAD
% -------------------------------------------------

\draw[hidden] (S)--(H);

% -------------------------------------------------
% Видимые рёбра основания
% -------------------------------------------------

\draw[visible] (A)--(B)--(C);

% -------------------------------------------------
% Видимые боковые рёбра
% -------------------------------------------------

\draw[visible] (S)--(A);
\draw[visible] (S)--(B);
\draw[visible] (S)--(C);

\end{tikzpicture}
}
